% ======================================================================
% QIJ — WSOM+2026 submission draft, v4
% Abstract + Sec 1 = author-edited (banked via chat); Secs 2-5 = Claude
% draft aligned to the banked vocabulary, awaiting author pass.
%
% TERMINOLOGY REGISTRY (current)
%   X = {x_i in R^d}, i=1..N   — dataset, sample size N
%   \hat F                     — empirical distribution of the sample
%   psi                        — influence function (defined Sec 2.1)
%   RF                         — receptive field (defined Sec 2.1);
%                                never bare "cell" or "field"
%   W, w_j, j                  — prototype set / prototype / RF index
%   n_j, p_j = n_j/N, p        — RF counts, mass fractions, vector
%   T(p)                       — estimator as function of mass fractions
%   V_tot = V_btw + V_win      — macros \Vtotal \Vcell \Vwithin
%   em dashes                  — only where the author wrote them;
%                                otherwise ", explanation" format
%   \todo{..}                  — awaiting runs or author confirmation
%
% OPEN ITEMS FROM SEC-1 EDITING (flagged, easy to revert):
%   [A] Intro decomposition para: I inserted the 5-word clause
%       "QIJ's estimate comprises the first." (recommended, not yet
%       explicitly approved) — delete the sentence to revert.
%   [B] The 22%-coverage empirical sentence remains OUT of Sec 1
%       (author rewrite omitted it; restoration was offered, no reply).
% ======================================================================
\documentclass[runningheads]{llncs}

\usepackage[T1]{fontenc}
\usepackage{graphicx}
\usepackage{amsmath,amssymb}
\usepackage{booktabs}
\usepackage{xcolor}
\usepackage{mathtools}

% --- draft machinery (delete before submission) ---
\newcommand{\todo}[1]{\textcolor{red}{[TODO: #1]}}
\newcommand{\placeholderfig}[2]{%
  \fbox{\parbox[c][#1][c]{0.94\linewidth}{\centering\ttfamily\small #2}}}

% --- notation ---
\newcommand{\ifl}{\psi}
\newcommand{\Vcell}{V_{\mathrm{btw}}}
\newcommand{\Vwithin}{V_{\mathrm{win}}}
\newcommand{\Vtotal}{V_{\mathrm{tot}}}
\newcommand{\Fhat}{\hat F}

\begin{document}

\title{Vector Quantization for Uncertainty Quantification: The Quantized Infinitesimal Jackknife}

\titlerunning{Vector Quantization for Uncertainty Quantification}
\author{\todo{Author One}\inst{1}\orcidID{\todo{0000-0000-0000-0000}} \and
\todo{Author Two}\inst{2}}
\authorrunning{\todo{A. One and A. Two}}
\institute{\todo{Institute One, City, Country}
\email{\todo{email@example.edu}}
\and
\todo{Institute Two, City, Country}}

\maketitle

\begin{abstract}
Nonparametric confidence intervals (CIs) for a parameter estimate $T(X)$ are commonly obtained by the bootstrap: recompute $T(X_b)$ on $B$ resampled datasets $X_b$ and infer estimate uncertainty from the replicates. The bootstrap is the gold standard, but it costs $B$ evaluations of the estimator, prohibitive when a single fit is expensive, e.g.\ when $T$ must be found by numerical optimization rather than in closed form. We show that a fitted vector quantizer offers a much cheaper route to the same intervals. The Quantized Infinitesimal Jackknife (QIJ) uses the quantizer's $M$ receptive fields (with counts $n_j$) as the variables of the uncertainty calculation: the estimator is re-expressed as a function of the counts, and QIJ measures how $T(X)$ responds when each $n_j$ is perturbed, which is the natural impact of resampling on a VQ. This requires only $M{+}1$ evaluations of $T(X)$ in total, vs.\ $B \gg M{+}1$ in the standard bootstrap. We infer the variance and skewness of the estimate by propagating the counts' exactly known (multinomial) sampling fluctuations through the measured responses; both feed our VQ-enabled CIs, and no resampling occurs anywhere. On two synthetic distributions (Pareto; multivariate $t$ in ten dimensions) and two astronomical datasets (the stellar initial mass function; the galaxy fundamental plane), QIJ matches the bootstrap's coverage and interval widths at a small fraction of its computation.
\keywords{vector quantization \and uncertainty quantification \and confidence intervals \and influence functions \and bootstrap}
\end{abstract}


% ======================================================================
\section{Introduction}\label{sec:intro}
% ======================================================================

Vector quantization (VQ) compresses a dataset $X = \{x_i \in \mathbb{R}^d\}_{i=1}^N$ into $M \ll N$ prototypes $W = \{w_j\}_{j=1}^M$ to facilitate signal encoding, clustering, density representation \cite{graf2000}, and, through self-organizing maps and their descendants, topology-preserving visualizations. In this work, we assign VQ a different task in the statistical inference domain: constructing confidence intervals for an arbitrary estimator computed from $X$.

The bootstrap \cite{efron1979} is the gold standard for nonparametric interval estimation: recompute the estimator $T$ on $B$ resampled datasets, with $B$ in the hundreds to thousands, and infer CIs from the replicated sampling distribution of $T(X)$. When a single evaluation of $T$ is expensive, e.g., the result of an optimizer-based fit, a profile likelihood, or an ensemble, those $B$ evaluations dominate the analysis. VQs represent large data well; it is natural to assume data prototypes (learned by the VQ) could stand in for the $N$ sample points during resampling. This works well for point estimates, but not for interval estimation. That is, in terms of uncertainty quantification, the quantizer prototypes are not just a smaller dataset from which to resample. Prototype representation sharpens signal and reduces sampling noise, but it is, exactly, this sampling noise that is needed for uncertainty quantification. We can, however, recover a notion of sampling variability by studying how resampling variance interacts with a VQ partition: there is variation in \emph{how many} points land in each region, and also in \emph{which} points represent each region. The VQ objective minimizes the second part by construction (if the VQ is trained properly); in this work we exploit multinomial models of the first to obtain CIs.

To achieve this, we move away from explicit resampling and instead turn to perturbation theory. Under bootstrap resampling, the counts $(n_1,\dots,n_M)$ of data in each region fluctuate with \emph{exactly known} (multinomial) covariance --- no simulation is required to learn it. The only unknown is how sensitive the estimator is to each changing count, and that is a derivative: perturb one region's share of the data, re-evaluate $T$, and record the response. Our Quantized Infinitesimal Jackknife (QIJ) does exactly this, via $M{+}1$ deterministic evaluations of $T$ in total, and propagates the counts' known covariance through these responses (the delta method), yielding the variance and skewness of $T$'s sampling distribution. Hence, we produce VQ-based confidence intervals, with no explicit resampling anywhere in the pipeline (Sect.~\ref{sec:method}).

These VQ-based confidence intervals are based on the observation that we can decompose an estimator $T$'s total variance into two parts. The first captures the estimator variance that would occur from distributing a resampling of points into the existing partition (i.e., Voronoi regions would obtain new counts), while the second describes estimator impact from within-region point variation. QIJ's estimate comprises the first. Crucially, the total variance (sum of these parts) is not dependent on a specific VQ partition; that is, decomposition under a different $M$ returns the same total estimator variance. We use this as a stringent correctness check in Sect.~\ref{sec:results}.

\paragraph{Related work.} The infinitesimal jackknife originates with Jaeckel \cite{jaeckel1972}; Efron \cite{efron1982} connects it to the bootstrap. Wager et al.\ \cite{wager2014}, similarly motivated by an estimator too expensive to resample, apply it to random forests, though their derivative is taken with respect to individual observation weights rather than regionally, as we do here. The Bag of Little Bootstraps \cite{kleiner2014} and $m$-out-of-$n$ subsampling \cite{bickel1997} reduce the size of each resample but still resample, and their approximation error is asymptotic and unobserved; QIJ never resamples, and its approximation error is a reported quantity. The influence-function machinery we rely on is classical robust statistics \cite{hampel1974,stefanski2002}.


% ======================================================================
\section{The Quantized Infinitesimal Jackknife}\label{sec:method}
% ======================================================================

\subsection{Setup: Estimators as Functions of Receptive Field Counts}\label{sec:setup}

This section develops the perturbation calculus previewed above. Let $X = \{x_i\}_{i=1}^N \stackrel{\mathrm{iid}}{\sim} F$ and let $T(X)$ be a function of $X$ to estimate unknown paramter $\theta$. $T$'s influence function, denoted $\ifl$ throughout, is
\begin{equation}
\ifl(x) \;=\; \lim_{\varepsilon \to 0}\frac{T\bigl((1-\varepsilon)F + \varepsilon\,\delta_x\bigr) - T(F)}{\varepsilon},
\end{equation}
which is the derivative of $T$ in the direction of a point mass at $x$ (i.e., it describes how much the estimate moves per unit of probability placed at $x$). The classical linearization $T(X) - \theta \approx \tfrac1N\sum_i \ifl(x_i)$ <<if this is clasical we need to cite it!>> underlies both the jackknife and the bootstrap, and gives $\mathrm{Var}\bigl(T(X)\bigr) \approx \tfrac1N\,\mathbb{E}[\ifl^2]$.

To begin, we fit a vector quantizer with $M$ prototypes $W = \{w_j\}$ to the data. Each prototype's Voronoi region is its \emph{receptive field} (RF); RF $j$ contains $n_j$ data points, with mass fraction $p_j = n_j/N$ and $p = (p_1,\dots,p_M)$. Any estimator that depends on the data through RF-level statistics can be written as a function $T(\cdot)$ of the mass fractions, such that $T(p) \coloneq T(\Fhat)$. We define the \emph{RF influence} as the derivative of $T$ toward RF $j$:
\begin{equation}\label{eq:Uj}
U_j \;=\; \frac{d}{d\varepsilon}\,T\bigl((1-\varepsilon)\,p + \varepsilon\, e_j\bigr)\Big|_{\varepsilon = 0}
\;=\; \frac{1}{n_j}\sum_{i \in \mathrm{RF}_j} \ifl(x_i),
\end{equation}
where $e_j \in \mathbb{R}^M$ is the $j$-th standard basis vector of the \emph{weight space}, i.e., the mass vector that places all weight on RF $j$ (note $e_j$ is unrelated to the data space $\mathbb{R}^d$). The perturbation raises RF $j$'s total weight by $\varepsilon$ at the expense of the others, reweighting every point inside RF $j$ identically, never individually; the response can therefore depend on the points of RF $j$ only through their average influence on $T$, which is the second equality. The RF influences are computable analytically when $\ifl$ has a closed form (or through the estimating-equation calculus of $M$-estimation \cite{stefanski2002}), and numerically otherwise, by a delete-one-RF secant: re-evaluate $T$ with RF $j$'s mass removed and the remainder renormalized, and divide the change by the mass moved. Either way, the full set $\{U_j\}$ costs $M{+}1$ deterministic evaluations of $T$, one at the observed $p$ and one per perturbed RF, and no randomness enters at any point.

The counts $(n_1,\dots,n_M)$ are multinomial, so their covariance is known exactly: $\mathrm{Cov}(p) = [\mathrm{diag}(p) - pp^{\top}]/N$. Propagating this known fluctuation through the derivatives (the delta method) gives the variance estimate
\begin{equation}\label{eq:Vcell}
\Vcell \;=\; \frac1N \sum_{j=1}^{M} p_j\, U_j^{2},
\end{equation}
using $\sum_j p_j U_j = 0$ to first order. The same quantities supply the third moment of the influence distribution, from which the skewness-based corrections of the BCa interval \cite{efron1987} follow with no additional evaluations. In short: the quantizer reduces the data to $M$ counts whose fluctuations are known; QIJ measures the estimator's sensitivity to each count; known fluctuation $\times$ measured sensitivity $=$ sampling distribution.

The counts obtained under resampling are jointly multinomial: redistributing $N$ draws into the fixed partition gives $Cov(n_j, n_k) = N(p_j \delta_{jk} - p_j p_k)$, with no simulation required to learn it. Thus, the corresponding mass fractions (the arguments of $T(\cdot)$) fluctuate with covariance $[ \text{diag}(p) - pp^T ] / N$. Propagating this known fluctuation through the derivatives via the delta method gives the variance estimate $$\Vcell \;=\; \frac1N \sum_{j=1}^{M} p_j\, U_j^{2}, \tag{3}$$ using $\sum_j p_j U_u = 0$ to first order. The same quantities supply the third moment of the influence distribution, from which the skewness-based corrections of the BCa interval \cite{efron1987} follow with no additional evaluations. In short: the quantizer reduces the data to $M$ counts whose fluctuations under resampling are known. QIJ measures the estimator's sensitivity to each count: known fluctuation × measured sensitivity gives us a sampling distribution.

\subsection{The Variance Decomposition}\label{sec:identity}

In this section, we study how compression to $M$ RFs allows decomposition of estimator variance.  With $U_j$ as above and $\mathrm{Var}_j(\ifl) = \tfrac{1}{n_j}\sum_{i \in \mathrm{RF}_j}(\ifl(x_i) - U_j)^2$ the spread of influence among the points inside RF $j$,
\begin{equation}\label{eq:identity}
\underbrace{\frac1N\sum_{i=1}^{N} \ifl(x_i)^{2}}_{N\,\Vtotal}
\;=\;
\underbrace{\sum_{j=1}^{M} p_j\, U_j^{2}}_{N\,\Vcell}
\;+\;
\underbrace{\sum_{j=1}^{M} p_j\, \mathrm{Var}_j(\ifl)}_{N\,\Vwithin}.
\end{equation}
This is the between/within variance decomposition applied to the influence values, and it holds algebraically on any dataset <<and any estimator?>> and any partition. 

We have two remarks about this decomposition. First, the left side contains no partition, while both right-side terms depend on it strongly: thus, the left-hand side is independent of $M$, while the right-hand side components change (but offset each other) with varying $M$. Second, $\Vwithin$ is computable from the variance of the per-point influence values $\psi_i$ of points within $RF_i$. Thus, comparing $\Vwithin$ to the total(via some proportionality threshold) tells a user whether $M$ was large enough to allow the VQ to resolve their estimand well. 

\subsection{Estimator Limitations}\label{sec:scope}

We note here certain characteristics of estimators for which QIJ is not applicable. Namely, from the decomposition above, we must have a measurable $V_{win}$ influence variance on $T$ from within each receptive field. This is not always the case; e.g., when estimating a quantile, it is very likely that all $\psi_i$ in RF_j have the same influence, and the above decomposition does not hold. Likewise, QIJ estimates estimator variance, and does not correct bias. Finally, we note that the multinomial covariance powering \eqref{eq:Vcell} assumes iid sampling; survey weights, selection corrections, and dependent data are outside the present framework.


% ======================================================================
\section{Experimental Data}\label{sec:data}
% ======================================================================

We evaluate on four data-generating processes: two synthetic families in which every quantity is analytically available, and two drawn from astronomy, chosen to test the claims where they are least convenient, with heavy tails, high dimension, vector-valued estimands, and an expensive optimizer-based fit. Throughout, $S$ Monte Carlo draws of size $N$ are taken from each DGP; QIJ and a standard nonparametric bootstrap ($B = \todo{B}$) run on identical draws, paired by seed \todo{confirm $S$, $N$, $B$, $M$ grids}. The astronomical cases use an \emph{empirical population} protocol: a large catalogue is treated as the population, draws are sampled from it with replacement, and the target $\theta_{\mathrm{true}}$, the estimand computed on the full catalogue, is known \emph{without error}; coverage statements on real data require no appeal to an assumed parametric truth.

\paragraph{Pareto.} Heavy-tailed and one-dimensional, with closed-form influence functions for both estimands (the shape parameter's maximum-likelihood estimator, and a tail probability), so every stage of the pipeline can be checked analytically.

\paragraph{Multivariate $t$.} Ten-dimensional and elliptically symmetric, with the degrees of freedom $\nu$ estimated by profile maximum likelihood, testing that the construction survives high dimension and a likelihood-based estimator.

\paragraph{Galaxy fundamental plane.} Early-type galaxies satisfy a tight empirical relation among velocity dispersion, surface brightness, and size \cite{djorgovski1987,dressler1987}. We use an SDSS sample of $79{,}000$ galaxies \cite{bernardi2003} as the empirical population; the estimand is the vector of plane coefficients together with the intrinsic scatter, giving real data and a vector-valued estimand.

\paragraph{Stellar initial mass function.} The IMF, the mass distribution of stars at formation \cite{salpeter1955,kroupa2001,chabrier2003}, is estimated from a catalogue of $\approx 20{,}000$ \todo{confirm} stellar masses formed in the STARFORGE simulations \cite{grudic2021}, treated as the empirical population. The estimand of record is the high-mass power-law slope, obtained by an optimizer-based fit \todo{one clause on the fitted form}. A single evaluation of $T$ is expensive enough that the bootstrap's $B$ evaluations dominate wall-clock time, making this the DGP on which the computational claim is measured rather than merely counted.


% ======================================================================
\section{Results}\label{sec:results}
% ======================================================================

\subsection{The Decomposition Is Exact, and the Intervals Are Right}\label{sec:res-valid}

Figure~\ref{fig:invariance} shows the identity \eqref{eq:identity} on a fixed dataset as $M$ sweeps \todo{range}: $\Vcell$ rises as the partition refines, $\Vwithin$ falls, and their sum is flat to \todo{precision} throughout, agreeing with the standard bootstrap's variance band. No parameter was fit to produce the flat line: the components come from different RF statistics, and their sum has no freedom to be constant unless the implementation is correct. The same figure shows directly what compression omitted, with $\Vwithin$ sitting \todo{orders of magnitude} below the total at operating $M$.

\begin{figure}[t]
\centering
\placeholderfig{4.2cm}{F3 --- decomposition invariance:\\ log--log, $x = M$; three series ($\Vcell$, $\Vwithin$, $\Vtotal$) per DGP panel;\\ horizontal band = bootstrap $V$. Fixed dataset, $M$ varying.}
\caption{The decomposition \eqref{eq:identity} is exact. For a fixed draw from each DGP, the between-RF term $\Vcell$ (rising) and within-RF term $\Vwithin$ (falling) trade off as the partition refines, while their sum (flat line) is constant in $M$ to \todo{precision} and matches the standard bootstrap's variance (shaded band).}
\label{fig:invariance}
\end{figure}

Correctness of the variance transfers to the intervals. Across DGPs and estimands, the paired ratio $\log(V_{\mathrm{QIJ}}/V_{\mathrm{boot}})$ lies within the $\pm 5\%$ materiality band (Table~\ref{tab:vratio}); offsets of order $1/N$ are expected, since QIJ omits higher-order terms and the bootstrap carries its own $O(1/N)$ bias. Figure~\ref{fig:coverage} shows both methods on the coverage diagonal within Monte Carlo error, with widths at matched coverage agreeing to \todo{X}\%. One second-order effect favors QIJ: its variance estimate, deterministic given the data, is more stable across draws than the bootstrap's (spread of $\log \hat V$ reduced \todo{X}\% after correcting the bootstrap arm for its finite-$B$ Monte Carlo floor).

\begin{table}[t]
\centering
\caption{QIJ vs.\ paired bootstrap variance: mean $\log(V_{\mathrm{QIJ}}/V_{\mathrm{boot}})$ with 95\% CIs, per DGP and estimand. Materiality band $\pm 0.05$. \todo{Fill from results table.}}
\label{tab:vratio}
\begin{tabular}{@{}llr@{}}
\toprule
DGP & Estimand & $\log$ ratio [CI] \\
\midrule
Pareto & shape & \todo{~} \\
Pareto & tail prob. & \todo{~} \\
Multivariate $t$ & $\nu$ & \todo{~} \\
Fund.\ plane & coefficients, scatter & \todo{~} \\
IMF & slope & \todo{~} \\
\bottomrule
\end{tabular}
\end{table}

\begin{figure}[t]
\centering
\placeholderfig{4.0cm}{F2 --- calibration: empirical vs nominal coverage,\\ one line per method, $y{=}x$ reference, MC-error band;\\ panel per DGP.}
\caption{Interval calibration. Empirical coverage against nominal level for QIJ (BCa intervals from $\{p_j, U_j\}$) and the standard bootstrap, with the $\sqrt{p(1-p)/S}$ Monte Carlo band. Both methods sit on the diagonal across the level grid; widths at matched coverage agree to \todo{X\%}.}
\label{fig:coverage}
\end{figure}

\subsection{The Same Interval at a Fraction of the Compute}\label{sec:res-compute}

Figure~\ref{fig:compute} reports the computational claim on the one DGP where it is measured rather than counted: the IMF slope, whose optimizer-based fit makes each evaluation of $T$ expensive. Sweeping the bootstrap over $B$ (by truncating stored replicates) traces its accuracy--cost frontier; QIJ sits at $M{+}1$ evaluations per $M$. At matched interval quality, QIJ requires \todo{X}$\times$ fewer evaluations of $T$ (\todo{wall-clock: $a$ s vs.\ $b$ s per draw}). Both methods estimate the same population quantity; QIJ reaches it deterministically, with the omitted term of \eqref{eq:identity} reported alongside.

\begin{figure}[t]
\centering
\placeholderfig{3.8cm}{F7 --- compute frontier (IMF):\\ $x$ = evaluations of $T$ (log); $y$ = interval error;\\ bootstrap swept over $B$; QIJ one point per $M$.}
\caption{Interval quality against computation for the IMF slope. The bootstrap frontier is traced by truncating stored replicates to $B' \le B$; QIJ (one point per $M$) attains the same interval quality at $M{+}1$ deterministic evaluations, \todo{X}$\times$ fewer evaluations of the optimizer-based estimator at matched quality.}
\label{fig:compute}
\end{figure}


% ======================================================================
\section{Conclusions}\label{sec:conclusions}
% ======================================================================

A fitted vector quantizer supports nonparametric interval estimation. The route is not resampling through the partition, which draws its randomness from precisely the variability the quantizer minimized, but perturbation across it: the RF counts fluctuate with exactly known covariance, the estimator's sensitivity to each count costs one evaluation, and $M{+}1$ deterministic evaluations therefore replace $B$ resamples. An exact identity keeps the accounting honest: the within-RF term that QIJ's perturbations cannot detect is computed rather than assumed, and the partition-invariance of the total provides a built-in correctness test. On four testbeds, including two real astronomical problems with empirically known truth, QIJ delivers nominal coverage at bootstrap-matched widths, at a small fraction of the computation.

The identity invites questions beyond parity that we develop in companion work: what the omitted term itself measures (its decay with $M$ turns out to carry the effective dimension of the data), and how a quantizer should be fitted when the goal is inference rather than distortion. Extending the multinomial engine to survey weights and dependent data, and shrinkage across RFs using prototype-adjacency structure, are further steps toward inference on real catalogues.

% \begin{credits}
% \subsubsection{\ackname} \todo{funding + acknowledgments}
% \subsubsection{\discintname} The authors have no competing interests. \todo{confirm}
% \end{credits}

\bibliographystyle{splncs04}
\bibliography{refs}

\end{document}
